\documentclass[journal,article,submit,pdftex,moreauthors]{Definitions/mdpi}

\firstpage{1}
\makeatletter
\setcounter{page}{\@firstpage}
\makeatother
\pubvolume{1}
\issuenum{1}
\articlenumber{0}
\pubyear{2026}
\copyrightyear{2026}
\datereceived{ }
\daterevised{ }
\dateaccepted{ }
\datepublished{ }

% ----- AIRAS: values realized by update_record; fallbacks until then -----
\InputIfFileExists{values.tex}{}{}
\providecommand{\airasval}[1]{\textbf{??airasval:\detokenize{#1}??}}
\providecommand{\unverified}[1]{#1}
\usepackage{listings}
\lstset{basicstyle=\ttfamily\small, breaklines=true, columns=fullflexible, keepspaces=true}

\Title{Preregistered Formal Verification of Two Summation Identities in Lean 4 through the AIRAS Integrity Flow}

\Author{AIRAS Agent $^{1}$}
\AuthorNames{AIRAS Agent}
\address{%
$^{1}$ \quad auto-res2 / lean-summation-identities-3 (GitHub)}
\corres{Correspondence: see the repository}

\abstract{AIRAS verifies experimental claims by preregistering them, executing the experiment in an isolated environment, and admitting into the paper only numbers that are re-derived from the recorded run outputs. This paper applies the same flow to a theoretical claim. Two elementary summation identities, the Gauss sum $2\sum_{i<n+1} i = n(n+1)$ and the sum of the first $n$ odd numbers $\sum_{i<n}(2i+1) = n^2$, are declared as claims whose verifier is Lean 4 with a pinned mathlib. Each claim names the module, the declaration and the exact type the declaration must have, and this type is frozen in the record before any proof is written. The proofs are then built on GitHub Actions, where a fixed report tool records the built type, the axioms the proof depends on, and the build's errors; the record's verdict is derived from that report alone. We preregister the two claims here and report, after the runs, whether each was supported or inconclusive, whether the built type matched the frozen statement, and which axioms were used.}

\keyword{Lean 4; mathlib; formal verification; preregistration; reproducibility; research integrity}

\begin{document}

\section{Introduction}

Research automation systems such as AIRAS~\cite{airas2025} make an experimental result trustworthy by separating three things: what the researcher declares in advance, what the computation actually produced, and what the paper reports. The declaration is frozen in version control before any run, the run executes on a platform the author cannot edit, and every number in the paper is re-derived from the run's stored outputs when the paper is checked. A claim that was reworded after the fact, or a number that was typed rather than measured, fails that check.

Theoretical claims have so far been outside this flow. A theorem is verified not by measuring a metric but by a proof checker, and the natural checker for present-day formalized mathematics is Lean~4 with its mathematical library mathlib~\cite{lean4,mathlib}. This paper asks whether the same discipline transfers: can a theorem be preregistered as a claim, proved afterwards, checked in isolation, and reported only through the checker's output?

We answer with a deliberately small case. The two identities below are textbook material, and their formal proofs are a few lines each. The point is not the mathematics but the flow: the \emph{statement} of each theorem is frozen before the proof exists, the proof is built where the author cannot intervene, and the verdict is derived mechanically from the build.

\section{Related Work}

Formal mathematics in Lean has grown into a large shared library, mathlib, and recent work uses it as training and evaluation data for machine-learning provers: Herald translates the mathlib corpus into natural language~\cite{gao-2025-herald}, Lean Workbook formalizes natural-language problems into Lean statements~\cite{ying-2024-lean}, and miniCTX evaluates provers on theorems whose proofs depend on repository context~\cite{hu-2025-minictx}. Draft, Sketch, and Prove guides an automated prover with informal proof sketches~\cite{jiang-2023-draft}, and POETRY proves theorems level by level, deferring subgoals with placeholders~\cite{wang-2024-proving}. All of these treat the proof checker as the arbiter of correctness; none of them place the \emph{statement} under a preregistration discipline, which is the gap this paper addresses.

Closed forms for finite sums have a classical algorithmic treatment~\cite{karr-1981-summation}; the identities we prove are the simplest instances of that theory.

\section{Background}

\paragraph{The AIRAS record.} A research repository holds a machine-readable record of hypotheses, claims, designs and runs. A claim carries a verifier. For an experimental claim the verifier is the execution platform and the claim states a criterion on a metric; for a theoretical claim the verifier is Lean, and the claim names a module, a declaration, and the declaration's type as Lean prints it. The record is append-only in git: a claim is never edited, only superseded by a later entry with the same identifier, so every version stays readable.

\paragraph{Lean claims.} A run of a Lean claim builds the module and then invokes a fixed report tool that loads the built environment, pretty-prints the declaration's type, collects the axioms the declaration depends on (the placeholder axiom \texttt{sorryAx} among them), and writes them with the build's errors and warnings to a report. The record's verdict is \emph{supported} when the report has no errors, the printed type equals the frozen statement, and every axiom is in the claim's allowed set; otherwise it is \emph{inconclusive}. Lean never refutes: a proof that did not go through shows nothing either way.

\section{Method}

The two identities are formalized over the natural numbers with mathlib's finite-sum notation. The proofs are by induction on $n$, unfolding the sum over $\mathrm{range}\,(n+1)$ with \texttt{Finset.sum\_range\_succ} and closing the arithmetic with the \texttt{ring} tactic. Each identity is one declaration in its own module under \texttt{lean/Airas/}. The Lean toolchain is pinned to \texttt{leanprover/lean4:v4.33.1} and mathlib to its \texttt{v4.33.1} tag; both pins are part of each claim's verifier, and a report built with a different toolchain or mathlib revision is an error.

\section{Hypothesis and Claims}

The hypothesis and its claims are rendered from the record; the list below is the record's, not a transcription, and it is regenerated and compared at every verification.

% AUTO-GENERATED by the update_record tool. DO NOT EDIT.
% verify_paper_values regenerates this file from record.json and the
% run outputs and fails on any difference, so manual edits always surface.
\noindent\textbf{H1.} Two elementary summation identities (the Gauss sum and the sum of the first n odd numbers) can be formally proved in Lean 4 with mathlib by induction, depending only on Lean's standard axioms, and that verification can run through the same preregistration, isolated-execution and provenance flow that AIRAS uses for experiments.
\begin{enumerate}
\item[\textbf{C1}] The Gauss sum, twice the sum of i over i below n+1 equals n(n+1), is proved in Lean 4 without sorry for every natural number n, and its proof depends on no axiom outside the standard three.
  \emph{Rationale:} The first instance of H1's claim that elementary summation identities are formally provable from the standard axioms alone. It is the most basic arithmetic-progression identity and should close with induction and Finset.sum\_range\_succ only.
  \emph{Verdict:} pending.
\item[\textbf{C2}] The sum of the first n odd numbers equals n squared, is proved in Lean 4 without sorry for every natural number n, and its proof depends on no axiom outside the standard three.
  \emph{Rationale:} The second instance of H1. It is an identity independent of C1, and closing it with the same induction pattern is evidence that the method does not depend on one particular identity.
  \emph{Verdict:} pending.
\end{enumerate}
\noindent\emph{Assumed, so that the claims together imply H1:}
\begin{itemize}
\item The statements of C1 and C2 faithfully formalize the intended mathematical content of each identity (faithfulness of a formalization is not something Lean checks).
\item The Lean 4 v4.33.1 kernel, mathlib v4.33.1 and the report tool (airas-report) are correct, and the axiom set returned by collectAxioms is complete.
\item The GitHub Actions run executes a descendant of the commit the record declared (the premise the provenance cross-check rests on).
\end{itemize}


Each claim's frozen statement is the exact type Lean must print for the declaration, recorded in \texttt{.research/record.json} under the run's \texttt{params.statement} and compared verbatim (up to whitespace) with the built declaration's type. In mathematical notation the two statements are:
\begin{itemize}
\item C1, run \texttt{gauss-sum}, module \texttt{Airas.GaussSum}, declaration \texttt{gauss\_sum\_mul\_two}: for every $n \in \mathbb{N}$, $2 \sum_{i \in \mathrm{range}(n+1)} i = n(n+1)$.
\item C2, run \texttt{odd-sum}, module \texttt{Airas.OddSum}, declaration \texttt{odd\_sum\_eq\_sq}: for every $n \in \mathbb{N}$, $\sum_{i \in \mathrm{range}(n)} (2i+1) = n^2$.
\end{itemize}

There is no numeric criterion and no predicted interval: the falsification line of a Lean claim is the build itself. A claim is inconclusive if the proof fails to build, uses \texttt{sorry}, is built for a type other than the frozen one, or depends on an axiom outside the allowed set. We expect both proofs to depend on \texttt{propext} at most, and possibly on no axiom at all.

\section{Experimental Setup}

Each run is dispatched to GitHub Actions on an \texttt{ubuntu-latest} runner (x86\_64, no GPU) at the \emph{full} stage, which requires a sorry-free proof. The repository's single entry point, \texttt{make run}, reads the run's kind from its configuration file, builds the named module with \texttt{lake build}, and writes the report with \texttt{lake exe airas-report}. The report is uploaded as the workflow artifact and imported into the repository together with a provenance manifest naming the workflow run, the commit it executed and the hash of every imported file. The record is then realized from the imported report, and the verification gate re-derives every value and byte-compares the imported files against the artifact.

% ===== airas: post-experiment sections, filled once runs exist =====
\section{Results}

Both claims were realized from the imported reports and both are \emph{supported}; the verdicts appear in the claim list of Section~5, which is regenerated from the record at every verification. For each run the report shows an empty error list, no warnings (so no \texttt{sorry}), and \texttt{statement\_matches} true: the declaration was built for exactly the type that was preregistered.

The statements were deliberately frozen in spellings other than the one Lean prints --- C1 with the arrow binder and \texttt{Finset.sum} applied to a lambda, C2 with the binder written without parentheses --- and the report tool, which elaborates the frozen text in the module's environment and compares it with the built type as a term, accepted both. A string comparison with the printed form would have rejected both.

The two proofs depend on the same three axioms, \texttt{propext}, \texttt{Classical.choice} and \texttt{Quot.sound} --- the standard axioms of Lean's core library, and exactly the allowed set declared for both claims. The preregistration expected \texttt{propext} at most; the difference is discussed below.

Each run executed on GitHub Actions in the repository's Docker image (elan with the pinned toolchain \texttt{leanprover/lean4:v4.33.1} and mathlib \unverified{\texttt{0df444a3}}) on an \texttt{ubuntu-latest} runner, on commit \unverified{\texttt{a420df7d}}, a descendant of the freeze commit \unverified{\texttt{9fab2ca3}}. The provenance manifest names the workflow runs (\unverified{34941959631} for \texttt{gauss-sum}, \unverified{34941963311} for \texttt{odd-sum}), the commit each executed, the dispatch stage (\texttt{full}) and the SHA-256 of every imported file; the verification gate byte-compares the imported reports against the runs' artifacts. Values quoted in this paragraph are transcribed from that manifest and marked as unverified, because the record has no declared value for them.

\section{Discussion}

\paragraph{The flow transfers.} Nothing in the preregistration, execution and verification machinery had to know that the claims were theorems rather than experiments. The claim declared a verifier, a module, a declaration and a frozen type; the same dispatch, status and import tools ran the proof and brought its report back; the same record realization derived the verdict; the same gate checked the record, the rendered claim list and the provenance. The one Lean-specific piece is the fixed report tool that prints the type, compares it with the frozen statement and collects the axioms, and it is deliberately not written by the author of the proof.

\paragraph{What the frozen statement buys.} The frozen statement is compared as a term, up to the names of bound variables, so the record may spell the type any way Lean accepts; what it cannot do is name a different theorem. Definitional unfolding is not used in the comparison: a proof of a statement that merely computes to the declared one (\(n = n\) for \(n + 0 = n\)) is not the declared theorem and does not match. The comparison runs inside the report tool, so it is only as trustworthy as that tool's copy in the run; a gate that rebuilds the proof with its own copy of the tool is the remaining step.

\paragraph{Axioms.} The expectation of \texttt{propext} at most was wrong: the tactics used (\texttt{simp}, \texttt{ring}) and mathlib's lemma \texttt{Finset.sum\_range\_succ} route through classical reasoning and quotient soundness. The proofs are no less valid for it --- the three axioms are the standard ones and were declared as allowed --- but it shows that the axiom set is worth predicting with the actual tactic vocabulary in mind, and that a claim wanting a constructive proof must say so in \texttt{allowed\_axioms}.

\paragraph{What a first trial found.} This is the second run of the same study. The first (auto-res2/lean-summation-identities-2) reached the same verdicts but exposed defects around the pipeline rather than in it: a stale Python lockfile made the Docker image fail to build and the workflow silently fell back to a runner without a Lean toolchain; the journal template's EPS logo could not be converted in a build without shell escape; the claim list did not escape LaTeX's special characters; the paper gate lacked the token to consult GitHub Actions' provenance; and the statement was compared as a printed string. All were fixed in the tooling before this run, which passed each of those points on the first attempt. None of the defects could have produced a wrong verdict --- a failed run has no report --- but they show that the execution environment and the rendering path are part of a claim's reproducibility.

\paragraph{Limitations.} Lean does not refute: a claim whose proof fails is inconclusive, not false, so the hypothesis could only have been supported or left open. The gate does not yet rebuild the proof itself; it checks that the imported report is the run's and that the record matches it, so the chain of trust runs through the workflow run and its artifact rather than through an independent re-derivation. The identities are elementary and the proofs short; the exercise tests the flow, not the difficulty of formalization. The faithfulness of each formal statement to its informal identity remains an assumption, as declared.

\section{Conclusions}

Two summation identities were preregistered as Lean claims in spellings of their statements that Lean itself does not print, proved after the freeze, built in an isolated GitHub Actions run, and reported only through the record realized from the run's report. Both claims are supported, with the frozen statements recognized as the built types by term comparison and the axioms within the declared set. The AIRAS integrity flow, built for experiments, carried a theoretical claim end to end; the residual work is a gate that re-derives the report by rebuilding the proof.

% ===== airas: end of post-experiment sections =====

\vspace{6pt}

\authorcontributions{The whole flow, from the hypothesis to this manuscript, was carried out by an AI agent under the AIRAS integrity rules; the proofs were checked by Lean on GitHub Actions.}
\funding{This research received no external funding.}
\dataavailability{All resources used in this study are openly available at \url{https://github.com/auto-res2/lean-summation-identities-3}.}
\acknowledgments{In this study, we automatically carried out a series of research processes---from hypothesis formulation to paper writing---using generative AI.}
\conflictsofinterest{The authors declare no conflicts of interest.}

\begin{adjustwidth}{-\extralength}{0cm}
\bibliographystyle{plainnat}
\bibliography{references}
\PublishersNote{}
\end{adjustwidth}
\end{document}
